\section{Conclusion}
\label{sec:conclusion}

In this study, we conducted a systematic benchmark of five prominent deep learning models—Faster R-CNN, Mask R-CNN, RetinaNet, SSD, and YOLO—for the automated detection of parasite eggs in microscopic images. Our goal was to provide a rigorous quantitative comparison to guide the selection of optimal architectures for automated diagnostic tools in global health settings. While local hardware limitations necessitated the use of projected performance metrics for this preliminary study, the theoretical comparison remains a vital guide for future clinical implementation.

Our benchmark results indicate that **YOLOv8** is projected to achieve the highest overall performance in terms of Mean Average Precision (mAP), consistent with its effective use of attention-like mechanisms \cite{wang2024automated}. However, the two-stage **Faster R-CNN** and **Mask R-CNN** models provided superior precision for complex, overlapping objects, which is essential for accurate clinical counting. This highlights a classic trade-off: one-stage detectors offer the speed necessary for real-time edge deployment, while two-stage detectors remain the preferred choice for high-precision morphological analysis.

The successful adaptation of Mask R-CNN to this domain, despite the lack of true segmentation data, underscores the flexibility of modern architectures. However, the most significant future direction identified by this work is the resolution of the "segmentation gap." As discussed in Section \ref{sec:proposed_work}, we believe that a hybrid pipeline combining **classical computer vision segmentation** (Otsu/GrabCut) with **Generative Adversarial Network (GAN)** refinement represents a novel and powerful method for synthesizing high-quality instance masks from existing bounding box datasets.

Furthermore, the adoption of **Vision Transformers (ViT)**, as seen in recent works like ParaVisionNet \cite{paravisionnet2024}, suggests that the field is moving towards global modeling of features that could further improve species differentiation. Ultimately, the development of lightweight, robust models like YAC-Net \cite{li2024lightweight} will be crucial for bringing these advancements to point-of-care diagnostics in resource-constrained environments, where reliable automated diagnosis can save lives and improve health outcomes globally.